% VI45-DETAILED-CHALLENGE-SOLUTIONS
% VI45-DETAILED-CHALLENGE-01
\subsection*{Challenge VI.45.1}
\begin{solution}
For \(y^2=x^3\), the first blow-up gives \(x^2(u^2-x)\).  The strict transform is smooth but tangent to the first exceptional divisor, so embedded normal crossings require another blow-up at that tangency point.  Each later blow-up decreases the self-intersection of every curve through its center by one and creates a new \((-1)\)-curve.  Continuing until the strict transform meets the exceptional configuration transversely gives a chain dual graph recording the cusp resolution.
\end{solution}

% VI45-DETAILED-CHALLENGE-02
\subsection*{Challenge VI.45.2}
\begin{solution}
Compose quadratic transformations so that an exceptional curve of one map lands at a base point of the next.  After the first proper-point blow-up, the residual base condition appears as a point on an exceptional divisor, i.e. an infinitely near point.  Blow it up and include the new exceptional class in the Picard basis.  The resulting lattice action is computed exactly as in the standard case but now must record proximity between successive exceptional divisors.
\end{solution}

% VI45-DETAILED-CHALLENGE-03
\subsection*{Challenge VI.45.3}
\begin{solution}
On the blow-up of three noncollinear points, the \(E_i\) and the three classes \(H-E_i-E_j\) all have square \(-1\) and anticanonical degree one.  Solving
\[
D^2=-1,\qquad (-K)\cdot D=1
\]
together with effectivity inequalities for three general points leaves precisely these six classes.  Hence they are the six \((-1)\)-curves of the degree-six del Pezzo surface.
\end{solution}

% VI45-DETAILED-CHALLENGE-04
\subsection*{Challenge VI.45.4}
\begin{solution}
For \(\alpha=H-E_1-E_2-E_3\), \(\alpha^2=-2\).  The reflection
\[
r_\alpha(D)=D+(D\cdot\alpha)\alpha
\]
reproduces the four Cremona basis formulas.  Because \(\alpha^2=-2\), it is involutive and preserves the intersection form.  Moreover \(K\cdot\alpha=0\), so the canonical class is fixed; the fixed subspace is \(\alpha^\perp\).
\end{solution}

% VI45-DETAILED-CHALLENGE-05
\subsection*{Challenge VI.45.5}
\begin{solution}
Algorithm: find the common zero scheme of a primitive homogeneous triple, blow up each proper base point in charts, divide the pulled-back triple by its common exceptional factor, inspect the exceptional divisor for residual common zeros, and recurse on those infinitely near points.  Record one exceptional class per blow-up and the moving-system class after each division.  When no common zero remains, the class defines the resolved morphism and its intersection data.
\end{solution}

